\section{Motivación}\label{cap:motivacion}

La principal motivación de este proyecto surge de la dificultad que tienen tanto los ciudadanos como los visitantes a la hora de buscar y obtener información clara y actualizada durante los días de la Semana Santa. Aunque a día de hoy existan programas oficiales, carteles informativos o cuentas oficiales en redes sociales, estos medios no siempre consiguen dar a conocer en tiempo real la situación de las procesiones, ya sea por retrasos o cambios de última hora debidos a factores meteorológicos o a problemas dentro de la propia cofradía.

Debido a esto, los asistentes, en muchas ocasiones, deben desplazarse a los distintos puntos del recorrido sin tener la certeza de que la procesión no haya pasado ya, cuánto tiempo de espera se estima hasta que pase por ese punto o incluso si llegará a pasar debido a cancelaciones. Esto genera desorganización y aglomeraciones innecesarias, además de una experiencia menos satisfactoria para los observadores.

Gracias a estas observaciones, se ve clara la necesidad de contar con un medio común capaz de organizar y difundir la información necesaria para que la experiencia sea lo más clara y accesible posible para cualquier persona.

Como respuesta a estas limitaciones, se plantea la idea de una aplicación móvil que permita la centralización de la información sobre el estado de las procesiones de forma actualizada, además de incluir otras funciones como la visualización de los recorridos y la geolocalización en tiempo real de las mismas. De esta forma, se mejora no solo la experiencia de los usuarios, sino que también se facilita la gestión de la información durante la Semana Santa.

Además, el desarrollo de esta aplicación permite no solo aplicar los conocimientos adquiridos durante el grado sobre análisis de requisitos, diseño de software, desarrollo de aplicaciones móviles y gestión de proyectos, sino también ampliarlos mediante su aplicación en un proyecto real, enfrentándose a problemas prácticos como la integración de distintas tecnologías, el manejo de la geolocalización en tiempo real o la gestión de múltiples tipos de usuarios dentro de un mismo sistema.
